\section{Exact Schur transport of the integral generator endpoint}
\label{sec:plethysm-endpoint-transport}

Retain the coefficient ring, original coefficient and original endpoint
factorization of Section~\ref{sec:integral-coefficient-conductor}:
\begin{gather*}
 O=\mathbb Z[1/2],\quad R=O[q,q^{-1}],\quad t=q-1,\quad
 r=q-q^{-1},\quad h=[7]_q[3]_q,\\
 C=q^{10}-q^4-q^{-4}+q^{-10}=r^2h.
\end{gather*}
The ordered rank-four modules and maps are
\[
 E_C\xrightarrow{H}B_R\xrightarrow{V}N_R,\qquad
 H=\operatorname{diag}(1,1,1,h),\quad
 V=\operatorname{diag}(1,1,1,r^2),\quad D_C=VH.
\]
The fourth coordinate of $E_C$ is the twisted source top line,
and that of $N_R$ is the bottom line, with the three specified identity
summands retained. Thus $D_C$ is the stabilized actual divided-square
endpoint, not a new full quantum-group representation. The displayed
bases identify each of the three modules with $R^4$ as an $R$-module;
we keep their names to distinguish domains and codomains.

All Schur functors in this section are the integral, column-exterior to
row-symmetric Schur functors described in
Theorem~\ref{thm:plethysm-integral-basis}. No factorial or hook
length is inverted. We apply these functors to ordinary module maps.
They retain the stated $\mu_2^2$ characters and the semilinear bar
involution, by functoriality. They also retain the explicitly diagonal
Cartan action and zero lowering actions on the stabilized endpoint
modules. This assertion does not apply an ordinary symmetric-group
action to a generic full quantum tensor power.

\subsection{The exact diagonal map for every partition}

For a partition $\lambda$, let $\operatorname{Tab}_d(\lambda)$ be its
tableaux with entries in $\{1,\ldots,d\}$ weakly increasing along rows
and strictly increasing down columns. The empty shape has one empty
tableau. A shape with more than $d$ rows has no such tableaux and its
Schur module is zero. For $T\in\operatorname{Tab}_4(\lambda)$ write
$w(T)$ for the number of entries equal to $4$. Put
\[
 d_\lambda=|\operatorname{Tab}_4(\lambda)|,\qquad
 m_\lambda(k)=|\{T:w(T)=k\}|,
 \qquad P_\lambda(z)=\sum_{k\ge0}m_\lambda(k)z^k.
\]
These definitions are finite and effective; they do not presume a
plethysm decomposition or a positivity conjecture.

\begin{theorem}\label{thm:plethysm-endpoint-diagonal}
In the corresponding tableau bases of the three modules, the maps
$S_\lambda(H)$, $S_\lambda(V)$ and $S_\lambda(D_C)$ have diagonal entries
$h^{w(T)}$, $r^{2w(T)}$ and $C^{w(T)}$, respectively. Their composition is exactly
$S_\lambda(D_C)=S_\lambda(V)S_\lambda(H)$. Each map is injective over
$R$, and
\[
 \operatorname{coker}S_\lambda(D_C)
 =\bigoplus_{k\ge1}\big(R/(C^k)\big)^{m_\lambda(k)}.
\]
Writing $\mathcal L_C,\mathcal L_r$ for their actual image lattices
inside $S_\lambda(N_R)$ gives the exact sequence
\begin{equation}\label{eq:plethysm-conductor-sequence}
 0\longrightarrow\bigoplus_{k\ge1}(R/(h^k))^{m_\lambda(k)}
 \xrightarrow{\ i\ }
 \bigoplus_{k\ge1}(R/(C^k))^{m_\lambda(k)}
 \xrightarrow{\ \pi\ }
 \bigoplus_{k\ge1}(R/(r^{2k}))^{m_\lambda(k)}
 \longrightarrow0,
\end{equation}
where $i(\bar f_T)=\overline{r^{2w(T)}f_T}$ and
$\pi(\bar f_T)=\bar f_T$. The three terms are, in that order,
$\mathcal L_r/\mathcal L_C$, $S_\lambda(N_R)/\mathcal L_C$,
and $S_\lambda(N_R)/\mathcal L_r$.
\end{theorem}

\begin{proof}
Each term in the column alternation followed by row multiplication
has exactly the content of $T$. A diagonal map therefore multiplies
the tableau vector by the product of its diagonal entries, giving
the three stated scalars. The integral tableau basis ensures these
are matrices over $R$ itself. Naturality, or multiplication of their
individual diagonal entries, gives the composition. Nonzero diagonal
entries are non-zero-divisors in the domain $R$, proving injectivity.
The cokernel of multiplication by a scalar $a$ on $R$ is $R/(a)$;
the coordinates with $w=0$ are identity maps and contribute zero.
For each $k>0$, multiplication by $r^{2k}$ sends $R/(h^k)$ into
$R/(C^k)$ since $r^{2k}h^k=C^k$. If its image is zero, the equality
$r^{2k}f=C^kg$ cancels to $f=h^kg$. Projection onto $R/(r^{2k})$
is onto and its kernel consists exactly of these multiples of
$r^{2k}$. Taking the finite direct sum proves every map and exactness
assertion, including the image-lattice interpretations.
\end{proof}

Pad $\lambda$ with zero rows to length four. The exact smallest and
largest exponents, for a nonzero Schur module, are
\begin{equation}\label{eq:plethysm-single-floors}
 \min_T w(T)=\lambda_4,\qquad \max_Tw(T)=\lambda_1.
\end{equation}
Every column of height four must contain $4$, so the lower bound is
$\lambda_4$. Fill row $i$ constantly with $i$ to attain it. At most
one $4$ occurs per column, giving the upper bound $\lambda_1$.
In a column of height $a$, place $5-a,\ldots,4$ in order. As column
heights weakly decrease from left to right, these entries weakly
increase along each row. This is a tableau attaining the upper bound.

The complete multiplicities also have an explicit smaller formula.
Deleting the entries $4$ leaves a shape $\mu$ with at most three rows
and
\[
 \lambda_1\ge\mu_1\ge\lambda_2\ge\mu_2\ge
 \lambda_3\ge\mu_3\ge\lambda_4.
\]
Conversely each tableau of that shape on $1,2,3$ extends uniquely by
filling $\lambda/\mu$ with $4$. Indeed the displayed inequalities say
precisely that the added boxes lie at row ends and no two share a
column; the row and column conditions then follow. Thus
\begin{equation}\label{eq:plethysm-branching}
 m_\lambda(k)=
 \sum_{\substack{\lambda_i\ge\mu_i\ge\lambda_{i+1}\ (1\le i\le3)\\
                  |\lambda|-|\mu|=k}}
 \frac{(\mu_1-\mu_2+1)(\mu_2-\mu_3+1)(\mu_1-\mu_3+2)}2.
\end{equation}
The displayed summand is an integer count, not a division in $R$.
To prove it directly, delete the $3$ entries next. The remaining
two-row shape $(a,b)$ satisfies
$\mu_2\le a\le\mu_1$ and $\mu_3\le b\le\mu_2$.
Deleting its entries $2$ leaves a one-row shape of any length from
$b$ to $a$, giving $a-b+1$ possibilities. Hence its total is
\[
 \sum_{a=\mu_2}^{\mu_1}\sum_{b=\mu_3}^{\mu_2}(a-b+1)
 =\frac{(\mu_1-\mu_2+1)(\mu_2-\mu_3+1)(\mu_1-\mu_3+2)}2,
\]
by the finite arithmetic-progression sums. This proves
\eqref{eq:plethysm-branching} without a dimension formula assumption.

\subsection{Nested Schur maps and the precise plethysm character}

Fix another partition $\mu$ with $d_\mu>0$, enumerate its integral
tableau basis $T_1,\ldots,T_{d_\mu}$ so that the weights
$w_i=w(T_i)$ weakly increase, and choose any fixed order within a
tie. A tableau $U\in\operatorname{Tab}_{d_\mu}(\lambda)$ indexes an
integral basis vector of $S_\lambda(S_\mu(E_C))$. Its exponent is
\[
 W(U)=\sum_{\text{boxes }b\text{ of }\lambda}w_{U(b)},\qquad
 M_{\lambda,\mu}(k)=|\{U:W(U)=k\}|.
\]
Apply Theorem~\ref{thm:plethysm-endpoint-diagonal} to a diagonal map
whose entries now are $C^{w_1},\ldots,C^{w_{d_\mu}}$. It proves, over
the unchanged ring $R$, the exact matrices
\begin{equation}\label{eq:plethysm-endpoint-nested-diagonal}
 S_\lambda(S_\mu(D_C))=\operatorname{diag}_{U}(C^{W(U)}),
 \quad S_\lambda(S_\mu(H))=\operatorname{diag}_{U}(h^{W(U)}),
 \quad S_\lambda(S_\mu(V))=\operatorname{diag}_{U}(r^{2W(U)}).
\end{equation}
All injectivity, cokernel and conductor formulas above hold with
$m_\lambda(k)$ replaced by $M_{\lambda,\mu}(k)$. This is a direct
application of the proved basis and scalar calculation; it assumes
neither exactness of Schur functors on arbitrary short exact sequences
nor an integral direct-sum decomposition of a plethysm.

For completeness define the symmetric polynomial
$s_\mu(x_1,\ldots,x_4)=\sum_T\prod_{b}x_{T(b)}$ by its tableau
character. The multivariable diagonal action proved above shows that
the character of the nested module is obtained by evaluating
$s_\lambda$ on the list of monomials $\prod_bx_{T_i(b)}$, with each
tableau contributing its own entry, including repeated monomials.
This substitution is precisely $s_\lambda[s_\mu]$ in four variables.
Consequently the exact one-variable identity is
\begin{equation}\label{eq:plethysm-character-specialization}
 \sum_kM_{\lambda,\mu}(k)z^k
 =s_\lambda(z^{w_1},\ldots,z^{w_{d_\mu}})
 =s_\lambda[s_\mu](1,1,1,z).
\end{equation}
Each coefficient on this restriction is a proved nonnegative integer
count. This equation does not determine the decomposition into
irreducible $GL_4$ characters: it is the restriction of a character
along one specified one-parameter subgroup. In particular this
computation supplies a transport map and its invariant factors, not
a plethysm-multiplicity obstruction for a class variety.

If $\lambda$ has at most $d_\mu$ rows, the exact extreme exponents are
\begin{equation}\label{eq:plethysm-nested-floors}
 a_{\lambda,\mu}=\sum_i\lambda_iw_i,
 \qquad b_{\lambda,\mu}=\sum_i\lambda_iw_{d_\mu+1-i}.
\end{equation}
In a column of height $a$, the entry in its $i$th box is at least
$i$ and at most $d_\mu-a+i$. Summing the corresponding ordered
weights proves the lower and upper bounds. Filling row $i$ by $i$
attains the lower bound. Filling each height-$a$ column with
$d_\mu-a+1,\ldots,d_\mu$ attains the upper bound and is row weakly
increasing, for the same reason as in
\eqref{eq:plethysm-single-floors}. Grouping the lower and upper
column sums by the row lengths gives exactly
\eqref{eq:plethysm-nested-floors}.

Let $d_0=m_\mu(0)=|\operatorname{Tab}_3(\mu)|$. Weight zero in the
nested module occurs precisely when every entry labels an inner
tableau of weight zero. Therefore
\begin{equation}\label{eq:plethysm-zero-rank}
 M_{\lambda,\mu}(0)=|\operatorname{Tab}_{d_0}(\lambda)|.
\end{equation}
This also handles $d_0=0$: only the empty outer shape has a
nonzero rank. At $C=0$ the matrices are diagonal with an identity
block of this size and zero elsewhere, giving their exact rank,
kernel and cokernel over every base-change ring, not only over fields.

For example $S_{(2)}(R^4)$ has inner weights $0^6,1^3,2^1$.
An outer two-box row chooses an unordered pair with repetitions.
The five weight counts for $S_{(2)}(S_{(2)}(D_C))$ are
\[
 (21,18,12,3,1):\quad
 \binom72,\quad6\cdot3,\quad6\cdot1+\binom42,\quad3\cdot1,\quad1.
\]
Its diagonal therefore has $21$ identity entries, $18$ entries $C$,
$12$ entries $C^2$, $3$ entries $C^3$, and one entry $C^4$.
No irreducible decomposition was needed for this exact matrix.

There is also a uniform determinant formula. Let $F$ be either
$S_\lambda$ or $S_\lambda\circ S_\mu$, let $n$ be its polynomial
degree and $d$ its rank on $R^4$. Then
\begin{equation}\label{eq:plethysm-determinant}
 \det F(D_C)=C^{nd/4},\qquad
 \sum_k k\,\operatorname{rank}F(R^4)_k=nd/4\in\mathbb Z.
\end{equation}
To prove the divisibility and exponent, grade the integral tableau
basis by the four original contents. Every vector has total content
$n$. A permutation matrix exchanging any two original basis vectors
induces, by functoriality, an integral isomorphism exchanging the
corresponding weight spaces. The sums of the four content coordinates
over all basis vectors are thus equal integers, whose total is $nd$.
The fourth sum is the sum of the exponents in the diagonal matrix.
Taking its determinant proves the formula, including its positive sign
in the matched ordered bases.

\subsection{Integral specialization, residual primes and exact orders}

For a weight $k>0$, put $Q_{h,k}=R/(h^k)$,
$Q_{C,k}=R/(C^k)$ and $Q_{r,k}=R/(r^{2k})$.
The resolution $0\to R\xrightarrow{t}R\to O\to0$ computes every
derived specialization used here. The ideal $(t)$ is prime, and
$h(1)^k=21^k\ne0$ in the domain $O$. Hence
$\operatorname{Tor}_1^R(Q_{h,k},O)=0$: if $tf=h^kg$, reduction modulo
$t$ forces $g(1)=0$, and cancellation gives $f\in(h^k)$.
The two other first Tor groups are free of rank one over $O$,
with their exact generators $C^k/t$ and $r^{2k}/t$. Indeed
$tf=C^kg$ forces $f=(C^k/t)g$, and this representative vanishes
modulo $C^k$ exactly when $g\in(t)$; the same proof applies to
$r^{2k}$. All higher Tor groups vanish by the length-one resolution.
The specialized sequence for this weight is therefore
\begin{equation}\label{eq:plethysm-derived-21}
 0\longrightarrow O\xrightarrow{21^k}O
 \xrightarrow{\mathrm{mod}\ 21^k}O/(21^k)
 \xrightarrow{0}O\xrightarrow{1}O\longrightarrow0.
\end{equation}
Projection sends $C^k/t$ to $h^k$ times $r^{2k}/t$, giving the
first displayed scalar. Lift $r^{2k}/t$ into $Q_{C,k}$ and multiply
by $t$: the result $r^{2k}$ is the image of $1\in Q_{h,k}$.
Thus the connecting map has the stated positive sign and sends the
chosen generator to $1$. The next map multiplies by $r^{2k}$ and
vanishes after specialization; the last is the identity.
The entire Schur or nested Schur sequence is the direct sum of
these sequences with the exact weight multiplicities already proved.

The specialized map $F(D_C)_0$ has kernel and cokernel
$O^{d-m(0)}$ in the positive-weight coordinates. The induced map
between its kernel and that of $F(V)_0$ is
$\operatorname{diag}_{k>0}(21^k)$ with repetitions $m(k)$.
It is injective and its cokernel is
\[
 \bigoplus_{k>0}(O/(21^k))^{m(k)}
 \simeq\bigoplus_{k>0}
       \big((O/(3^k))\oplus(O/(7^k))\big)^{m(k)}.
\]
For the final isomorphism use the map taking a class to its two
residues. Since $3^k$ and $7^k$ are coprime, choose integers $a,b$
with $a3^k+b7^k=1$; an inverse sends $(x,y)$ to
$b7^kx+a3^ky$ modulo $21^k$. Its two compositions are the identity.
This proves the two residual prime modules without discarding the
original factor $21^k$.

It is essential to specify which valuation is being used.
At the height-one vertical prime $(\ell)$ of $R$, for $\ell=3$ or
$7$, the polynomial $C$ is a unit of $R_{(\ell)}$, since its Laurent
coefficients include $1$. Every transported map is consequently an
isomorphism there; its vertical valuations are zero. The ideal
$(\ell,t)$ gives a two-dimensional local ring and does not itself
specify a discrete valuation. The exceptional orders below instead
refer to the discrete valuation rings
\[
 D_0=\mathbb Q[q,q^{-1}]_{(t)},\qquad
 D_\ell=\mathbb F_\ell[q,q^{-1}]_{(t)}\quad(\ell\text{ odd}).
\]
Write $u=(q+1)/q$, so $r=tu$ and $u(1)=2$. The exact orders and
leading coefficients of the original $C$ are
\[
\begin{array}{c|c|c}
\text{coefficient field}&\operatorname{ord}_t C&
          t^{-\operatorname{ord}_tC}C\bmod t\\\hline
\mathbb Q&2&84\\
\mathbb F_3&4&1\\
\mathbb F_7&8&5\\
\mathbb F_\ell, \ell\ne2,3,7&2&84\bmod\ell.
\end{array}
\]
Here $C=t^2u^2h$ gives the characteristic-zero and ordinary-prime
rows. For an odd prime $\ell$, multiply the finite sum for
$[\ell]_q$ by $q^{\ell-1}$ and then by $q^2-1$. In characteristic
$\ell$ the result is $q^{2\ell}-1=(q^2-1)^\ell$.
Cancellation in the polynomial domain gives
$[\ell]_q=(q-q^{-1})^{\ell-1}$. Hence $C=r^4[7]_q$ in
characteristic $3$ and $C=r^8[3]_q$ in characteristic $7$.
Their leading residues are $2^4\cdot7=1$ in $\mathbb F_3$ and
$2^8\cdot3=5$ in $\mathbb F_7$, proving the table exactly.

Let $v$ be the relevant middle-column integer $2,4$, or $8$, and
$c$ the last-column residue. Raising the retained equality to the
$k$th power proves
\begin{equation}\label{eq:plethysm-smith-orders}
 C^k=t^{vk}\varepsilon_k,\quad\varepsilon_k\in D_\ell^\times,
 \quad\varepsilon_k\bmod t=c^k,\qquad
 \operatorname{coker}F(D_C)\simeq
 \bigoplus_{k>0}(D_\ell/(t^{vk}))^{m(k)}.
\end{equation}
This is an explicit Smith form: multiply each matrix row by
the inverse unit $\varepsilon_k^{-1}$, equivalently use target basis
vector $\varepsilon_k$ times the original vector, retaining that unit and its
specified leading residue as part of the original map. Ordered by
$k$, consecutive powers of $t$ divide one another. Its length is
$v\sum_kkm(k)$, and the exact least and largest matrix-entry orders are $v$ times
the exponents in \eqref{eq:plethysm-single-floors} or
\eqref{eq:plethysm-nested-floors}.
More precisely, the $j$th filtration quotient of the cokernel, for
$j\ge0$, has residue-field dimension
\[
 \dim\big(t^j\operatorname{coker}/t^{j+1}\operatorname{coker}\big)
 =\sum_{k:vk>j}m(k).
\]
For a single cyclic summand the classes $1,t,\ldots,t^{vk-1}$ are
a basis, proving both length and filtration formulas by direct sum.

The mixed-prime arithmetic can also be retained before reduction.
In $\mathbb Z_\ell[[t]]$, with $q=1+t$ and $\ell=3$ or $7$, define
\[
 E_\ell(t)=\Phi_\ell(1+t)
 =\frac{(1+t)^\ell-1}{t}
 =\sum_{j=0}^{\ell-1}\binom\ell{j+1}t^j.
\]
Here the quotient is the displayed polynomial identity. Its constant
term is $\ell$, its leading term is $t^{\ell-1}$, and each other
coefficient is divisible by $\ell$; the constant term is not
divisible by $\ell^2$. The exact quantum factor identity
\[
 [\ell]_q=q^{1-\ell}\Phi_\ell(q)\Phi_{2\ell}(q),\qquad
 \Phi_{2\ell}(q)=\sum_{j=0}^{\ell-1}(-q)^j,
\]
follows by multiplying the two geometric sums; $\Phi_{2\ell}(1)=1$.
The other quantum factor in $h$ has value $7$ when $\ell=3$ and
$3$ when $\ell=7$. Thus, writing that factor as $h_{\ne\ell}$,
the full equality is
\begin{equation}\label{eq:plethysm-mixed-local}
 C^k=t^{2k}E_\ell(t)^k U_\ell(t)^k,
 \quad U_\ell=u^2q^{1-\ell}\Phi_{2\ell}(q)h_{\ne\ell}
           \in\mathbb Z_\ell[[t]]^\times.
\end{equation}
The constant terms of $U_3,U_7$ are respectively $28$ and $12$,
which are units in the stated coefficient rings. This retains the
horizontal factor $t^{2k}$ and the separate arithmetic factor
$E_\ell^k$; it does not replace a two-dimensional local ring by an
unspecified valuation. Modulo $\ell$, $E_\ell=t^{\ell-1}$, recovering
exactly the orders $4k$ and $8k$.

There is an additional first Tor term if the conductor sequence is
first reduced to characteristic $3$ or $7$ and then specialized at
$t=0$. Over the corresponding Laurent polynomial ring put
$\bar h=h\bmod\ell$; it is nonzero and divisible by $t$.
The three first Tor generators are now
$\bar h^k/t$, $\bar C^k/t$, and $\bar r^{2k}/t$.
Multiplication by $\bar r^{2k}$ sends the first exactly to the
second; projection sends the second to $\bar h^k$ times the third,
which is zero on the residue field. The connecting map again sends
the third generator to $1$. Therefore, for every positive weight,
the full sequence over $\mathbb F_\ell$ is
\[
 0\longrightarrow\mathbb F_\ell\xrightarrow{1}\mathbb F_\ell
 \xrightarrow{0}\mathbb F_\ell\xrightarrow{1}\mathbb F_\ell
 \xrightarrow{0}\mathbb F_\ell\xrightarrow{1}\mathbb F_\ell
 \longrightarrow0.
\]
These descriptions follow from the same length-one resolution;
in particular the extra Tor term has been calculated, not omitted.

Finally, all exponents in the original integral maps are nonnegative.
A negative valuation does occur for the actual inverse after the
specified localization: $F(D_C)^{-1}$ has entries $C^{-k}$ over
$R[C^{-1}]$ and orders $-vk$ over the fraction field of $D_\ell$.
If a positive-weight coordinate occurs, this inverse cannot be an
$R$-linear inverse on the original lattices: an inverse image of its
unit target vector would require $C^{-k}\in R$, which is impossible
since $C$ vanishes at $q=1$. Its precise failure to be integral is
the cokernel already computed in
\eqref{eq:plethysm-smith-orders}. This supplies an exact pole/defect
correspondence without identifying a negative Laurent coefficient
with a negative plethysm multiplicity.

\subsection{Exact negative inertia at the retained circle character}

Retain the circle real structure $q^*=q^{-1}$, with conjugation on
coefficients, and the original exact character from the conic calculation:
\[
 q_0=\frac{\sqrt{14}+i\sqrt2}{4},\quad
 q_0^{-1}=\frac{\sqrt{14}-i\sqrt2}{4},\quad
 p_0=\frac{\sqrt{14}}2,\quad r_0=\frac{i}{\sqrt2},\quad
 a_0=\frac12,\quad s_0=\frac14.
\]
The coefficient field is $K_0=\mathbb Q(\sqrt{14},i\sqrt2)\subset\mathbb C$
with its displayed embedding and usual conjugation. The product of the
two displayed values for $q_0$ is $(14+2)/16=1$, so evaluation is a
homomorphism respecting the circle involution. To check every factor,
put $x_0=q_0^2+q_0^{-2}=p_0^2-2=3/2$. Then
\[
 [3]_{q_0}=x_0+1=\frac52,\qquad
 [7]_{q_0}=x_0^3+x_0^2-2x_0-1=\frac{13}{8}.
\]
The second identity expands $[7]_q$ as
$(q^6+q^{-6})+(q^4+q^{-4})+(q^2+q^{-2})+1$
and uses $q^6+q^{-6}=x^3-3x$, $q^4+q^{-4}=x^2-2$.
Consequently the full factorization retains the exact signs
\begin{equation}\label{eq:plethysm-circle-coefficient}
 r_0^2=-\frac12,\qquad h(q_0)=\frac{65}{16}>0,
 \qquad C(q_0)=-\frac{65}{32}.
\end{equation}

Let $F$ be either Schur functor already constructed, and identify its
source and target \emph{as based vector spaces} by the map taking
each source tableau basis vector to the equally labelled target
vector. This specifies the endomorphism $A_F$ representing
$F(D_C)(q_0)$ on the source; it does not infer an equality of the
original source and target modules. Give this source the positive
Hermitian form with its specified tableau basis orthonormal.
Because every diagonal entry of $A_F$ is the real number
$(-65/32)^k$, the endomorphism is selfadjoint for this form.
The Hermitian form that it defines has the complete formula
\[
 B_F(u,v)=\sum_T(-65/32)^{w(T)}\overline{u_T}v_T
\]
for a single Schur functor, and the same formula with $U,W(U)$ for
a nested one. In particular its positive, negative and zero indices
are exactly
\begin{equation}\label{eq:plethysm-inertia}
 n_+=\sum_{k\text{ even}}m(k),\qquad
 n_-=\sum_{k\text{ odd}}m(k),\qquad n_0=0.
\end{equation}
Indeed the even and odd coordinate spans give a direct orthogonal
sum of positive and negative definite spaces. A subspace of dimension
larger than the even-coordinate count intersects the odd-coordinate
space nontrivially by the dimension formula, so cannot be positive
definite. Interchanging even and odd proves maximality for the
negative index. Every diagonal coefficient is nonzero, so the
radical is zero. This proves the inertia statement over the specified
embedded field and after extension to $\mathbb C$.

The signature is thus the exact character evaluation
\[
 n_+-n_-=P_\lambda(-1)=s_\lambda(1,1,1,-1)
\]
for a single Schur functor, and
\[
 n_+-n_-=s_\lambda[s_\mu](1,1,1,-1)
 =s_\lambda(\underbrace{1,\ldots,1}_{m_\mu(\mathrm{even})},
             \underbrace{-1,\ldots,-1}_{m_\mu(\mathrm{odd})})
\]
for a nested one. The last equality follows by evaluating each
inner-tableau monomial at $(1,1,1,-1)$; it is $+1$ or $-1$ according
to its weight parity. The scalar $h(q_0)^k$ is positive for every
coordinate. Therefore $F(V)(q_0)$, whose corresponding coefficients
are $(-1/2)^k$, has exactly the same inertia. The actual conductor
factor $F(H)(q_0)$ is positive definite in these matched bases and
their matrices still multiply to $F(D_C)(q_0)$. This proves the
propagation of the negative directions through the actual endpoint
factorization, without discarding the positive residual factor.

For $S_{(2)}\circ S_{(2)}$ the previously computed profile gives
$n_+=21+12+1=34$, $n_-=18+3=21$, and signature $13$.
This also shows that the signature need not stay negative under
iteration even when some negative directions survive: the exact
parity-weight character, not an unsigned dimension, determines it.
No entry of $m(k)$ has become negative. These Hermitian forms are
specified forms on the endpoint Schur modules; no positivity form
of the nonstandard RH programme or spectral membership of $s_0$
has been identified with them.

\subsection{The original shape and an actual specialized source projector}

For the original shape $\nu=(7,5,3)$, formula
\eqref{eq:plethysm-branching} is the following explicit finite sum:
\begin{gather*}
 \sum_{a=5}^{7}\sum_{b=3}^{5}\sum_{c=0}^{3}
 \frac{(a-b+1)(b-c+1)(a-c+2)}2\,z^{15-a-b-c}\\
 =27+81z+162z^2+270z^3+300z^4+252z^5+126z^6+42z^7.
\end{gather*}
The exponents of $S_\nu(D_C)$, with multiplicities, are therefore
exactly the eight displayed coefficients. Their sum is $1260$ and
their weighted sum is $4725$. In characteristic zero at $q=1$,
the map has rank $27$, kernel and cokernel dimension $1233$, and
cokernel length $9450$ over $D_0$. The corresponding lengths over
$D_3,D_7$ are $18900$ and $37800$. These follow by multiplying
$4725$ by the proved orders $2,4,8$. They describe this particular
integral Schur transport and its stated base changes.
At the distinct circle character $q_0$ of
\eqref{eq:plethysm-circle-coefficient}, its inertia is
\begin{gather*}
 (n_+,n_-,n_0)
 =(27+162+300+126,\ 81+270+252+42,\ 0)=(615,645,0),\\
 \operatorname{signature}=-30.
\end{gather*}
This is a negative signature with $645$ individually specified
negative tableau directions, and is an immediate exact instance of
\eqref{eq:plethysm-inertia}.

There is also a proved map to the full original source, over its
already justified rational specialization. Put
$X=V\otimes W$ with the unchanged alphabet
\[
 x_1=v_1\otimes w_1,\quad x_2=v_1\otimes w_2,\quad
 x_3=v_2\otimes w_1,\quad x_4=v_2\otimes w_2.
\]
Section~\ref{sec:full-source-generators} constructs the isomorphism
$s_\nu:L_{\mathbb Q}\to Q_\nu(X)$ from the rational specialized
full source lattice to the exterior-column classical relation quotient,
and proves its inverse $j_\nu$. Here $Q_\nu(X)$ denotes that precise
quotient presentation, retaining its original column order. To compare
it with the integral Schur-image convention in this section, define
\[
 \eta_\nu:Q_\nu(X)\longrightarrow S_\nu(X),\qquad
 [\text{column exterior tensor of }T]\longmapsto b_T.
\]
Every classical adjacent column relation is the alternating shuffle
relation: the free letters form an alternating list of length $a+1$
across two columns with left height $a$, with the other letters fixed.
Its row-polynomial image vanishes by
Lemma~\ref{lem:plethysm-shuffle}; the unchanged other columns multiply
this zero polynomial. Thus $\eta_\nu$ is well defined. It is onto
because column tensors generate the Schur image. Its source has
dimension $1260$ by the proved source isomorphism, and its target has
dimension $1260$ by the explicit tableau count above. Its kernel is
therefore zero. Its inverse has a basis formula: send $b_T$ for each
semistandard $T$ to its column-tensor class. This right inverse is
also a left inverse by injectivity, proving both compositions.
Alternation, row multiplication and quotient maps are equivariant,
so both maps are equivariant as well. Put
\[
 \sigma_\nu=\eta_\nu s_\nu:L_{\mathbb Q}\longrightarrow S_\nu(X),
 \qquad\sigma_\nu^{-1}=j_\nu\eta_\nu^{-1}.
\]
Use the explicit coordinate isomorphisms
$\theta_E:E_C\otimes_R\mathbb Q\to X$ and
$\theta_N:N_R\otimes_R\mathbb Q\to X$ taking their $i$th displayed
basis vector to $x_i$. Since $C(1)=0$, their square is
\[
 \theta_N(D_C)_0=P_X\theta_E,\qquad
 P_X(x_i)=x_i\ (1\le i\le3),\quad P_X(x_4)=0.
\]
Applying $S_\nu$ and composing with the proved source maps gives the
actual endomorphism
\begin{equation}\label{eq:plethysm-original-source-projector}
\Pi_\nu=\sigma_\nu^{-1} S_\nu(P_X)\sigma_\nu:
 L_{\mathbb Q}\longrightarrow L_{\mathbb Q}.
\end{equation}
It is idempotent: $\sigma_\nu\sigma_\nu^{-1}=1$, functoriality, and $P_X^2=P_X$
give $\Pi_\nu^2=\Pi_\nu$. The Schur tableau formula shows its image
is $\sigma_\nu^{-1}S_\nu(\langle x_1,x_2,x_3\rangle)$ with its $27$-element
tableau basis, and its kernel is the span
of the $1233$ tableaux containing $x_4$, pulled back by
$\sigma_\nu^{-1}$. These are explicit inverse
descriptions of both subspaces, proving the ranks independently of
an arbitrary choice of $1260$ vectors in the source.

The original diagonal $GL(V)\times GL(W)$ torus preserves each
$x_i$, so it commutes with $P_X$. The proved equivariance of
$\sigma_\nu,\sigma_\nu^{-1}$ and Schur functoriality prove that $\Pi_\nu$ commutes
with this torus. No full raising/lowering equivariance is claimed:
already $F_Vx_2=x_4$, so $P_XF_Vx_2=0$ whereas
$F_VP_Xx_2=x_4$. The construction provides the exact coordinate
and source maps under which the endpoint specialization acts.
The rational source isomorphism is used only over $\mathbb Q$;
the integral torsion and the characteristic-$3,7$ calculations
above have not been transferred through an unproved integral or
modular source isomorphism.

The complete transport thus reaches ordinary Schur and nested Schur
modules, their exact conductor quotients, their residual prime
modules and their specified source specialization. It provides no
class-variety coordinate-ring obstruction, canonical-basis
Frobenius model, or Boolean complexity reduction by itself. Those
maps must be supplied by the continuing GCT construction; no
unproved complexity conclusion is a hypothesis in any result here.
